\section*{ID6332: What I Computed and Plotted: φ≈1.618034}
\paragraph{Metadata.}
PiID: 4090389064495 | RTEID: RTE:P24588.5622:T2.8153 | UUID: 2bc5db19-c58d-5054-ab2d-1b4f68c82b07

\paragraph{Canonical Statement.}
\#\#\# What I computed and plotted I numerically evaluated \textbackslash{}(I(\textbackslash{}ell)\textbackslash{}) for a range of \textbackslash{}(\textbackslash{}ell\textbackslash{}) values (including integer 1, integer 2, and the golden ratio \textbackslash{}(\textbackslash{}varphi\textbackslash{}approx1.618034\textbackslash{})). I plotted: - \textbackslash{}(|I(\textbackslash{}ell)|\textbackslash{}) (magnitude) vs. \textbackslash{}(\textbackslash{}ell\textbackslash{}), and - \textbackslash{}(\textbackslash{}arg I(\textbackslash{}ell)\textbackslash{}) (phase) vs. \textbackslash{}(\textbackslash{}ell\textbackslash{}).

\paragraph{Canonical Equation.}
\[
\varphi\approx1.618034
\]

\paragraph{Dictionary.}
Relation: φ≈1.618034

\paragraph{Encyclopedia.}
What I Computed and Plotted: φ≈1.618034 is a scaling / order-of-magnitude relation, written φ≈1.618034. It expresses a scaling or proportionality, capturing how the quantity grows with its parameters. It is built from φ. Within the Trawin Zero Point registry it serves as one element of the framework's mathematical narrative.
